\section{Data Preprocessing}
Before applying any machine learning method, it is essential to ensure that the dataset is clean, consistent,
and properly scaled. In this project, the preprocessing pipeline consists of three main steps:
\begin{enumerate}
	\item handling missing values,
	\item normalization of numeric features,
	\item splitting the data into training and testing sets.
\end{enumerate}

All preprocessing steps were implemented in Python using the \texttt{pandas} and \texttt{numpy} libraries.
The corresponding code is organized into modular functions in the Jupyter notebook in order to keep the
implementation clean and reusable.

\subsection{Handling Missing Values}
After loading the integrated heart disease dataset as a CSV file, we first inspected the data using
\texttt{df.info()} and \texttt{df.isna().sum()} to check for missing values and potential inconsistencies.
To ensure a complete dataset, we applied a simple but robust imputation strategy:

\begin{itemize}
	\item For \textbf{numeric} features (such as \texttt{age}, \texttt{trestbps}, \texttt{chol},
	\texttt{thalach}, and \texttt{oldpeak}), missing values were replaced with the
	\emph{median} of the corresponding column.
	\item For \textbf{categorical or binary} features (such as \texttt{sex}, \texttt{cp},
	\texttt{fbs}, \texttt{restecg}, \texttt{exang}, \texttt{slope}, \texttt{ca}, and \texttt{thal}),
	missing values were replaced with the \emph{mode} (most frequent value) in each column.
\end{itemize}

Median imputation is less sensitive to outliers than mean imputation and is therefore suitable
for medical measurements that may contain extreme values. Mode imputation for categorical
variables preserves the most common category and avoids introducing unrealistic new values.

The imputation was implemented in a dedicated function \texttt{impute\_missing\_values()}, which
takes the dataset and the lists of numeric and categorical columns as input and returns a cleaned
version of the data. After this step, we verified that the number of missing values in all columns
is zero.

\subsection{Normalization of Numeric Features}
Because the KNN classifier is a distance-based algorithm, it is crucial that all numeric features
are on a comparable scale. Otherwise, attributes with larger numeric ranges would dominate the
distance computation and bias the classification result.

In this project, we applied \textbf{Min--Max scaling} to all numeric features. For a numeric
attribute $x$, the normalized value $x'$ is computed as:
\begin{equation}
	x' = \frac{x - x_{\min}}{x_{\max} - x_{\min}},
\end{equation}
where $x_{\min}$ and $x_{\max}$ are the minimum and maximum values of that feature in the
\emph{training} set.

To avoid data leakage, the scaling parameters (minimum and maximum) were computed using only
the training data. The same parameters were then applied to transform the test data.
This was implemented in the function \texttt{min\_max\_scale\_train\_test()}, which returns
scaled versions of both the training and test sets.

If a numeric column happened to be constant (i.e., $x_{\max} = x_{\min}$), the normalized
values for that column were set to zero to avoid division by zero.

\subsection{Train--Test Split}
To evaluate the generalization performance of the KNN classifier, the dataset was split into
a training set and a testing set. The data were first randomly shuffled and then divided using
an \textbf{80\% / 20\%} split ratio:
\begin{itemize}
	\item 80\% of the instances were used for training the model,
	\item 20\% of the instances were held out as an independent test set.
\end{itemize}

The split was implemented in a helper function \texttt{train\_test\_split\_df()}, which takes
the full dataset and a \texttt{test\_size} parameter as input. A fixed random seed was used to
ensure reproducibility of the results.

After the split and scaling, the final matrices \texttt{X\_train}, \texttt{y\_train},
\texttt{X\_test}, and \texttt{y\_test} were obtained. These were used in the subsequent
experiments with the KNN classifier.